\documentclass[10pt]{scrartcl}
\usepackage[mdthm,colorsec,noauthor,noasy]{evan}
\usepackage{handout}
\title{Math 110 --- Week 1}
\subtitle{Solutions}
\author{}\date{}\ihead{}
\begin{document}
\maketitle

\begin{prob}\textbf{HW1 P1.}\par\idea{Each statement is about membership, so take one element and unfold the two definitions. For (c), ask which elements could be lost or gained.}
\begin{steps}
  \ii (a) Let $y\in f(f^{-1}(V))$. Then $y=f(x)$ for some $x\in f^{-1}(V)$, and $x\in f^{-1}(V)$ means $f(x)\in V$. So $y\in V$.
  \ii (b) Let $x\in U$. Then $f(x)\in f(U)$, which is the definition of $x\in f^{-1}(f(U))$.
  \ii (c) Strict in (a): $f\colon\RR\to\RR$, $f(x)=x^2$, $V=\{-1,1\}$: $f^{-1}(V)=\{\pm1\}$ and $f(f^{-1}(V))=\{1\}$. Strict in (b): same $f$, $U=\{1\}$: $f^{-1}(f(U))=\{\pm1\}$.
  \ii Equality in (a) for a given $V$ holds exactly when $V\subseteq f(X)$: an element $y\in V$ is in $f(f^{-1}(V))$ exactly when $y=f(x)$ for some $x$. It holds for every $V$ exactly when $f$ is surjective.
  \ii Equality in (b) for a given $U$ holds exactly when no point outside $U$ has the same value as a point of $U$. It holds for every $U$ exactly when $f$ is injective: if $f$ is injective and $f(x)=f(u)$ with $u\in U$ then $x=u\in U$; if $f(x)=f(x')$ with $x\ne x'$, take $U=\{x\}$.
\end{steps}
\end{prob}

\begin{prob}\textbf{HW1 P2.}\par\idea{Compute the square, then one more product.}
\begin{steps}
  \ii Let $\omega=\frac{-1+\sqrt3i}2$. $\omega^2=\frac{1-2\sqrt3i+3i^2}{4}=\frac{-2-2\sqrt3i}{4}=\frac{-1-\sqrt3i}{2}$.
  \ii $\omega^3=\omega^2\omega=\frac{(-1-\sqrt3i)(-1+\sqrt3i)}{4}=\frac{1-3i^2}{4}=\frac{4}{4}=1$.
\end{steps}
\end{prob}

\begin{prob}\textbf{HW1 P3.}\par\idea{If $a\ne0$ it has an inverse; multiplying by it isolates $b$. This needs $0\cdot x=0$, which is proved from the distributive law.}
\begin{steps}
  \ii $0\cdot x=(0+0)x=0\cdot x+0\cdot x$; adding $-(0\cdot x)$ gives $0\cdot x=0$.
  \ii If $ab=0$ and $a\ne0$: $b=(a^{-1}a)b=a^{-1}(ab)=a^{-1}\cdot0=0$.
\end{steps}
\end{prob}

\begin{prob}\textbf{HW1 P4.}\par\idea{Same as P3, with a vector in place of $b$.}
If $a\ne0$: $v=1\cdot v=(a^{-1}a)v=a^{-1}(av)=a^{-1}\cdot0=0$, using axiom (v) and Prop.~4(ii).
\end{prob}

\begin{prob}\textbf{HW1 P5.}\par\idea{The zero function decides one direction; linearity of the integral gives the other.}
\begin{steps}
  \ii If the set is a subspace it contains the zero function, whose integral is $0$. So $b=0$.
  \ii If $b=0$: the zero function is in the set. If $\int f=\int g=0$, then $\lambda f+g$ is continuous and $\int(\lambda f+g)=\lambda\int f+\int g=0$.
\end{steps}
\end{prob}

\begin{prob}\textbf{HW1 P6.}\par\idea{A subspace of a complex vector space must be closed under multiplication by complex scalars.}
No. $(1,0)\in\RR^2$ but $i\cdot(1,0)=(i,0)\notin\RR^2$. ($\RR^2$ is a subspace of $\CC^2$ regarded as a \emph{real} vector space.)
\end{prob}

\begin{prob}\textbf{HW1 P7.}\par\idea{Every cross term carries a factor $2$, and $2=0$ in $\FF_2$.}
\begin{steps}
  \ii $(x+1)^2=x^2+2x+1=x^2+1$.
  \ii $(x+1)^4=(x^2+1)^2=x^4+2x^2+1=x^4+1$.
  \ii $(x+1)^8=(x^4+1)^2=x^8+2x^4+1=x^8+1$.
  \ii Over $\RR$ the middle terms $2x$, $2x^2$, $2x^4$ do not vanish, and $(x+1)^8$ has nine nonzero terms with coefficients $1,8,28,56,70,56,28,8,1$. Over $\FF_2$ all the even coefficients are $0$.
\end{steps}
\end{prob}

\begin{prob}\textbf{R1.}\par\idea{Intersection: both conditions hold in each subspace. Union: closure under addition fails.}
\begin{steps}
  \ii Let $U_\alpha$ be subspaces. $0$ is in every $U_\alpha$. If $v,w$ are in every $U_\alpha$, so is $\lambda v+w$. So the intersection is a subspace.
  \ii The union need not be: the two axes in $\RR^2$ contain $(1,0)$ and $(0,1)$ but not their sum. See E1 for exactly when it is.
\end{steps}
\end{prob}

\begin{prob}\textbf{E1.}\par\idea{If neither contains the other, pick one vector in each that is not in the other and add them; the sum cannot be in either.}
\begin{steps}
  \ii If $U\subseteq W$ then $U\cup W=W$ is a subspace; likewise if $W\subseteq U$.
  \ii Suppose $U\cup W$ is a subspace and $U\not\subseteq W$. Pick $u\in U$ with $u\notin W$. Let $w\in W$ be arbitrary. Then $u+w\in U\cup W$.
  \ii If $u+w\in W$, then $u=(u+w)-w\in W$, which is false. So $u+w\in U$, and then $w=(u+w)-u\in U$. Hence $W\subseteq U$.
\end{steps}
\end{prob}

\begin{prob}\textbf{E2.}\par\idea{For two subspaces, direct means the intersection is $\{0\}$.}
\begin{steps}
  \ii $U\cap W=\{(0,y,0)\}\ne\{0\}$: not direct.
  \ii $U\cap X$: $(t,t,t)$ with third coordinate $0$ forces $t=0$. Direct. Similarly $W\cap X=\{0\}$ (first coordinate $0$): direct.
  \ii $U+X=\RR^3$: $(a,b,c)=(a-c,b-c,0)+(c,c,c)$. So $\RR^3=U\oplus X$.
\end{steps}
\end{prob}

\end{document}
